\Chapter{Odtwarzalność i pochodzenie}
\label{ch:repro}

\section{Wersje źródeł i dowodów}

Pełna rewizja kodu opisana w raporcie to
\begin{quote}\small
\nolinkurl{5ba376f904b58caf2072ce807cc822db7a90ccc6}.
\end{quote}
Odczytano pliki wersjonowane dokładnie w tej rewizji, niezależnie od
bieżącego drzewa roboczego. Kod poprawek T2 przed scaleniem ma SHA
\code{e4d484d}; producent kandydata danych ma SHA \code{e96ae05}.
Audyt T0 po poprawkach wykonano na \code{49a38c9}.
Identyfikatory pełnej długości, selektory liczb i SHA-256 plików dowodowych
zapisano poza narracją w \code{evidence-map.md} i
\code{evidence/provenance.json}. Plik \code{evidence/summary.json}
zawiera wyłącznie wybrane agregaty, bez surowych danych i prywatnych
dzienników.

Wynik 209 testów pochodzi z dziennika końcowych kontroli T2,
a pokrycie z \path{reports/T2-evidence/coverage-summary.json}.
Ustalenia T0 opisuje \path{docs/fastf1-findings.md} i dowód
\path{reports/T0-fix1-evidence/summary-after.json}.
Katalog \code{reports/} oznacza wydzielone archiwum dowodowe projektu,
częściowo niepubliczne. Skróty umożliwiają sprawdzenie tożsamości dowodu,
ale nie zastępują dostępu do niego. Publiczne źródła wskazano w bibliografii.

Nie powtarzano pobrań, audytu pełnego cache, treningu ani testowej
ewaluacji modeli na potrzeby napisania raportu. Włączenie kodu do
\code{main}, wynik testów i akceptacja danych są trzema różnymi stanami.

\section{Odtworzenie środowiska projektu}

Poniższe polecenia dotyczą repozytorium programu w opisanej rewizji.
Pierwsze przygotowuje środowisko; kolejne kontrole wykorzystują dane
testowe i blokadę sieci. Nie uruchamiają pobrania danych wyścigowych.

\begin{lstlisting}
uv sync --locked --extra train
uv run --offline ruff check .
uv run --offline ruff format --check .
uv run --offline pytest
uv run --offline pitstop --help
\end{lstlisting}

Osoba mająca właściwe, rozpakowane archiwum może sprawdzić je bez
cache i poświadczeń:
\begin{lstlisting}
uv run --offline --no-sync pitstop verify-data \
  --manifest data/raw/v1/manifest.json
\end{lstlisting}
Ścieżka musi wskazywać manifest wewnątrz wersjonowanego katalogu danych,
a nie jego oddzielną kopię. Uzyskanie kodu 0 potwierdza zakres kontroli
weryfikatora; nie zastępuje wymaganej akceptacji człowieka.

\section{Budowanie dokumentu}

Raport jest samodzielny. Plikiem głównym jest \code{main.tex},
kompilatorem pdfLaTeX, a bibliografię obsługuje BibTeX:
\begin{lstlisting}
latexmk -pdf -interaction=nonstopmode \
  -halt-on-error -outdir=build main.tex
\end{lstlisting}
Wynik to \code{build/main.pdf}. Nie potrzeba zewnętrznych danych,
fontów, sieci ani \code{shell-escape}. Instrukcje Overleaf i aktualizacji
zawiera anglojęzyczny \code{README.md}.
